\section{Encoded frames and small-patch rewrites}
\label{sec:peps_encoded_frames}
The statements below formalize Definition~6.1 and the first assertion of
Lemma~6.3 of \cite{OpenAI2026PolynomialPEPS}, September 24, 2026,
\texttt{05-frames.tex}, lines 11--97 and 182--252, the registers of a frame
with one register per site and per tag, together with the steps of the proof
of the second assertion listed below. The expansion of the branches into
allowed monomials on these registers is in
Section~\ref{sec:peps_small_rewrite_monomials}, and the second assertion with
its polynomial count is
Theorem~\ref{thm:peps_frame_small_rewrite_polynomial}.

Throughout, $\Lambda$ is a finite set of sites with positions in the plane,
each site carries a raw register $\C^q$ with $q\ge1$, and
$\mathcal H_\Lambda=\bigotimes_{x\in\Lambda}\C^q$. For $D\subseteq\Lambda$ and
an operator $X$ on $\mathcal H_D$, write $X_D\otimes\id$ for the operator
acting as $X$ on the sites of $D$ and as the identity elsewhere. Operators
supported on disjoint sets of sites commute. A vector $\Omega$ is fixed, and
$\lVert\cdot\rVert$ denotes the Euclidean norm of vectors and the operator
norm of matrices, rectangular ones included.

\begin{definition}[Square cylinder data]
    \label{def:peps_frame_square_patch}
    \lean{TNLean.PEPS.EncodedFrame.squareSample,
        TNLean.PEPS.EncodedFrame.SquarePatch,
        TNLean.PEPS.EncodedFrame.SquarePatch.inner,
        TNLean.PEPS.EncodedFrame.SquarePatch.outer,
        TNLean.PEPS.EncodedFrame.SquarePatch.Tag,
        TNLean.PEPS.EncodedFrame.SquarePatch.cyl,
        TNLean.PEPS.EncodedFrame.SquarePatch.proj,
        TNLean.PEPS.EncodedFrame.SquarePatch.proj_of_outer_eq_empty}
    \leanok
    For a center $c\in\R^2$ and a scale $h>0$, the inner and outer squares
    $Q^-(c,h)$ and $Q^+(c,h)$ are the closed sup-norm balls of radii $h$ and
    $2h$; their physical samples are the sites whose positions lie in them.
    Square cylinder data at $(c,h)$ consist of radii
    $h\le r_1\le\cdots\le r_n\le 2h$, the samples $D_j$ of the squares of
    radius $r_j$, and unit vectors $v_{j\ell}\in\mathcal H_{D_j}$ for
    $\ell\le d_j$, such that the cylinder projectors
    $\bigl(\ket{v_{j\ell}}\bra{v_{j\ell}}\bigr)_{D_j}\otimes\id$ have mutually
    orthogonal ranges. Their sum is
    \[
        P_{c,h}=\sum_{j=1}^{n}\sum_{\ell=1}^{d_j}
        \bigl(\ket{v_{j\ell}}\bra{v_{j\ell}}\bigr)_{D_j}\otimes\id .
    \]
    Every $D_j$ contains the inner sample and is contained in the outer
    sample. If the outer sample is empty, then $P_{c,h}=\id$, whatever the
    cylinder data. This is the projector of \texttt{05-frames.tex}, lines 26--48,
    with the cylinder data of \cite[Proposition~4.1]{OpenAI2026PolynomialPEPS};
    the accuracy $\lVert(\id-P_{c,h})\Omega\rVert\le\varepsilon$ enters as
    a hypothesis where it is used. The polynomial bound on $\sum_jd_j$ is not
    recorded: the count of branches below takes a uniform bound on the number
    of cylinder terms of each patch instead.
\end{definition}

\begin{definition}[Hole encoder]
    \label{def:peps_frame_hole_encoder}
    \lean{TNLean.PEPS.EncodedFrame.SquarePatch.branch,
        TNLean.PEPS.EncodedFrame.SquarePatch.encoder,
        TNLean.PEPS.EncodedFrame.stack,
        TNLean.PEPS.EncodedFrame.SquarePatch.card_tag_of_outer_eq_empty,
        TNLean.PEPS.EncodedFrame.SquarePatch.encoder_submatrix_of_outer_eq_empty}
    \leanok
    \uses{def:peps_frame_square_patch}
    Let $\mathcal J_{c,h}$ be a tag register with orthonormal basis
    $\{\ket{j,\ell}\}$ and let $\ket0$ be a fixed one-site unit vector. The
    hole encoder is
    \[
        K_{c,h}=\sum_{j,\ell}\ket{j,\ell}\otimes
        \bigl[\bigl(\ket0^{\otimes D_j}\bra{v_{j\ell}}\bigr)_{D_j}
            \otimes\id\bigr]
        \colon\mathcal H_\Lambda\to\mathcal J_{c,h}\otimes\mathcal H_\Lambda .
    \]
    It keeps a raw register at every site, inside the hole as well. If the
    outer sample is empty, the tag register is one-dimensional and $K_{c,h}$
    is the identity encoding, as in \texttt{05-frames.tex}, lines 67--69.
\end{definition}

\begin{lemma}[Encoders are contractions]
    \label{lem:peps_frame_encoder_contraction}
    \lean{TNLean.PEPS.EncodedFrame.SquarePatch.proj_isStarProjection,
        TNLean.PEPS.EncodedFrame.SquarePatch.encoder_conjTranspose_mul_self,
        TNLean.PEPS.EncodedFrame.SquarePatch.norm_encoder_le_one}
    \leanok
    \uses{def:peps_frame_hole_encoder}
    The operator $P_{c,h}$ is an orthogonal projection,
    $K_{c,h}^\dagger K_{c,h}=P_{c,h}$, and $\lVert K_{c,h}\rVert\le1$.
    This is \texttt{05-frames.tex}, equation (encoder-contraction),
    lines 61--63.
\end{lemma}

\begin{proof}\leanok
    If the outer sample is empty, $K_{c,h}$ is the identity encoding and all
    three claims are immediate. Otherwise, each cylinder projector is Hermitian
    and idempotent because $v_{j\ell}$
    is a unit vector, and products of distinct cylinder projectors vanish by
    orthogonality; hence $P_{c,h}$ is an orthogonal projection. Orthogonality
    of the tag basis gives
    $K_{c,h}^\dagger K_{c,h}=\sum_{j,\ell}
        (\ket{v_{j\ell}}\braket{0^{\otimes D_j}}{0^{\otimes D_j}}\bra{v_{j\ell}})_{D_j}
        \otimes\id=P_{c,h}$. Finally
    $\lVert K_{c,h}\rVert^2=\lVert K_{c,h}^\dagger K_{c,h}\rVert
        =\lVert P_{c,h}\rVert\le1$.
\end{proof}

\begin{definition}[Encoded frame]
    \label{def:peps_frame}
    \lean{TNLean.PEPS.EncodedFrame.Hole,
        TNLean.PEPS.EncodedFrame.Frame,
        TNLean.PEPS.EncodedFrame.Frame.raw,
        TNLean.PEPS.EncodedFrame.TagSpace,
        TNLean.PEPS.EncodedFrame.rawProd,
        TNLean.PEPS.EncodedFrame.frameEncoder,
        TNLean.PEPS.EncodedFrame.Frame.encoder,
        TNLean.PEPS.EncodedFrame.Frame.refVec}
    \leanok
    \uses{def:peps_frame_hole_encoder}
    An encoded frame $\mathcal F$ consists of a raw ownership assignment on
    one sheet, a list of holes $(c,h)$ with square cylinder data and pairwise
    disjoint outer samples, and a specified party holding each hole's tag.
    The order of the list is the fixed ordering of the appended tags. The
    encoding $K_{\mathcal F}$ is the product of the hole encoders, with the
    tags appended in that order, and the reference vector is
    $\Omega_{\mathcal F}=K_{\mathcal F}\Omega$. A frame without holes is a raw
    sheet. This is \cite[Definition~6.1]{OpenAI2026PolynomialPEPS},
    \texttt{05-frames.tex}, lines 71--82. The source requires the outer
    squares to be disjoint in the plane; only the disjointness of their
    physical samples, which follows, is used.
\end{definition}

\begin{lemma}[Disjoint encodings commute]
    \label{lem:peps_frame_commute}
    \lean{TNLean.PEPS.EncodedFrame.rawProd_commute,
        TNLean.PEPS.EncodedFrame.rawProd_append,
        TNLean.PEPS.EncodedFrame.tagAppendEquiv}
    \leanok
    \uses{def:peps_frame}
    For two lists of holes with disjoint outer footprints, the raw parts of
    their encodings commute on every pair of tag configurations. The tags of
    a concatenated list are the pairs of tags of the two parts, and on such a
    pair the raw part of the encoding of the concatenation is the product of
    the raw parts of the two encodings. Hence the encodings commute after
    the canonical identification of tag orderings
    (\texttt{05-frames.tex}, lines 84--85).
\end{lemma}

\begin{proof}\leanok
    The raw part on a tag configuration is a product of operators supported
    in the outer footprints of the holes of the list, so the two products
    are supported on disjoint sets of sites. The factorization follows by
    induction on the first list.
\end{proof}

\begin{lemma}[Frame Gram identity]
    \label{lem:peps_frame_gram}
    \lean{TNLean.PEPS.EncodedFrame.sum_rawProd_conjTranspose_mul_self,
        TNLean.PEPS.EncodedFrame.frameEncoder_conjTranspose_mul_self,
        TNLean.PEPS.EncodedFrame.projProd_isStarProjection,
        TNLean.PEPS.EncodedFrame.norm_frameEncoder_le_one}
    \leanok
    \uses{lem:peps_frame_encoder_contraction,lem:peps_frame_commute}
    For a frame with hole projectors $P_1,\dots,P_r$,
    $K_{\mathcal F}^\dagger K_{\mathcal F}=P_1\cdots P_r$, this product is an
    orthogonal projection, and $\lVert K_{\mathcal F}\rVert\le1$.
\end{lemma}

\begin{proof}\leanok
    By induction on the list. Write the raw part on the tags $(s,t)$ as
    $B_sR_t$, with $B_s$ from the first hole. Summing $R_t^\dagger
        B_s^\dagger B_sR_t$ over $s$ gives $R_t^\dagger P_1R_t$, and $P_1$
    commutes with $R_t^\dagger$ by disjointness of supports; the sum over $t$
    is then $P_1$ times the identity for the remaining holes. The projectors
    commute pairwise, so their product is an orthogonal projection of norm at
    most one.
\end{proof}

\begin{lemma}[Telescoping on a vector]
    \label{lem:peps_frame_telescoping}
    \lean{TNLean.PEPS.EncodedFrame.norm_act_list_prod_sub_le,
        TNLean.PEPS.EncodedFrame.norm_act_projProd_sub_le}
    \leanok
    For contractions $A_1,\dots,A_s$ and every vector $\psi$,
    \[
        \lVert(A_1\cdots A_s-\id)\psi\rVert
        \le\sum_{j=1}^s\lVert(A_j-\id)\psi\rVert .
    \]
    In particular, if each of $r$ projectors has error at most $\varepsilon$
    on $\psi$, so does their product up to $r\varepsilon$; no commutation is
    assumed (\texttt{05-frames.tex}, lines 237--248).
\end{lemma}

\begin{proof}\leanok
    Induction on $s$, from
    $(A_1R-\id)\psi=A_1(R-\id)\psi+(A_1-\id)\psi$ and $\lVert A_1\rVert\le1$.
\end{proof}

\begin{theorem}[Norm of the frame reference vector]
    \label{thm:peps_frame_norm}
    \lean{TNLean.PEPS.EncodedFrame.Frame.norm_refVec_eq,
        TNLean.PEPS.EncodedFrame.Frame.norm_refVec_le,
        TNLean.PEPS.EncodedFrame.Frame.one_sub_mul_le_norm_refVec}
    \leanok
    \uses{lem:peps_frame_gram,lem:peps_frame_telescoping}
    For a frame with $r$ holes, $\lVert\Omega_{\mathcal F}\rVert
        =\lVert P_1\cdots P_r\Omega\rVert$. If $\lVert\Omega\rVert=1$ and every
    hole projector satisfies $\lVert(\id-P_a)\Omega\rVert\le\varepsilon$, then
    \[
        1-r\varepsilon\le\lVert\Omega_{\mathcal F}\rVert\le1 .
    \]
    This is \texttt{05-frames.tex}, equation (frame-norm), lines 86--94.
\end{theorem}

\begin{proof}\leanok
    The two vectors have the same squared norm
    $\braket{\Omega}{P_1\cdots P_r\Omega}$, because the product of the
    projectors is an orthogonal projection. The upper bound is
    contractivity, and the lower bound is the reverse triangle inequality
    with the telescoping lemma.
\end{proof}

\begin{definition}[Registers of a frame]
    \label{def:peps_frame_registers}
    \lean{TNLean.PEPS.EncodedFrame.siteRegs,
        TNLean.PEPS.EncodedFrame.tagRegs,
        TNLean.PEPS.EncodedFrame.sites,
        TNLean.PEPS.EncodedFrame.nodup_sites,
        TNLean.PEPS.EncodedFrame.mem_sites,
        TNLean.PEPS.EncodedFrame.sitesEquiv,
        TNLean.PEPS.EncodedFrame.owner_siteRegs,
        TNLean.PEPS.EncodedFrame.layoutRegs,
        TNLean.PEPS.EncodedFrame.Frame.regs,
        TNLean.PEPS.EncodedFrame.Frame.map_owner_regs}
    \leanok
    \uses{def:peps_frame,def:peps_party_layout}
    The registers of a frame $\mathcal F$ are one register $\C^{\mathcal
            J_a}$ for each hole $a$, held by the party holding its tag, listed in
    the order of the holes, followed by one raw register $\C^q$ for each site
    $x\in\Lambda$, held by the raw owner of $x$, listed in a fixed order of
    the sites. This list contains every site exactly once, so taking the
    site at a position gives a bijection from its positions to $\Lambda$.
    More generally, if all sites in any list have raw owner $p$, every raw
    register of that list is owned by $p$. The tensor product of the frame's
    registers is its memory in the
    sense of Theorem~5.2 of \cite{OpenAI2026PolynomialPEPS}. This is the
    sheet of \texttt{05-frames.tex}, lines 13--20, with one raw register at
    each site and a specified owner for each register, together with the tags
    of Definition~6.1, lines 71--82.
\end{definition}

\begin{lemma}[Canonical coordinates of the registers]
    \label{lem:peps_frame_registers_coordinates}
    \lean{TNLean.PEPS.EncodedFrame.siteVec,
        TNLean.PEPS.EncodedFrame.siteVec_congr,
        TNLean.PEPS.EncodedFrame.siteIsoList,
        TNLean.PEPS.EncodedFrame.siteIsoList_siteVec,
        TNLean.PEPS.EncodedFrame.groupIso,
        TNLean.PEPS.EncodedFrame.groupIso_siteVec,
        TNLean.PEPS.EncodedFrame.tagVec,
        TNLean.PEPS.EncodedFrame.tagIso,
        TNLean.PEPS.EncodedFrame.tagIso_tagVec,
        TNLean.PEPS.EncodedFrame.layoutVec,
        TNLean.PEPS.EncodedFrame.layoutIso,
        TNLean.PEPS.EncodedFrame.layoutIso_layoutVec,
        TNLean.PEPS.EncodedFrame.layoutIso_symm_single,
        TNLean.PEPS.EncodedFrame.Frame.regIso}
    \leanok
    \uses{def:peps_frame_registers}
    The tensor product of the registers of a frame is identified, by a
    unitary, with $\C^{\mathcal J\times[q]^\Lambda}$, the space of the tags
    and the sheet on which the encoding $K_{\mathcal F}$ and its operators
    act: the product basis vector $\ket\tau\otimes\bigotimes_x\ket{c(x)}$ is
    sent to $\ket{\tau,c}$. More generally, the raw registers of a list of
    distinct sites enumerating a set $A$ are identified with $\C^{[q]^A}$ in
    the same way. These are the canonical site and sheet coordinates of
    \texttt{05-frames.tex}, lines 94--96, in which vectors are compared while
    the owners of their registers are specified separately.
    For any list $\ell$ of sites, including lists with repetitions, write
    $\ket c_\ell$ for its product basis vector. If $c(x)=c'(x)$ for every
    site in $\ell$, then $\ket c_\ell=\ket{c'}_\ell$.
\end{lemma}

\begin{proof}\leanok
    By induction on the list of registers, identifying $\C^\alpha\otimes
        \C^\beta$ with $\C^{\alpha\times\beta}$ at each step; product basis
    vectors are sent to standard basis vectors, and a list of distinct sites
    enumerating $A$ gives a bijection between its positions and $A$.
    Configurations agreeing on a list give the same basis vector at every
    position, hence the same product vector, even if sites repeat.
\end{proof}

\begin{definition}[Data and conditions of a small-patch rewrite]
    \label{def:peps_frame_rewrite_data}
    \lean{TNLean.PEPS.EncodedFrame.SmallPatchRewrite,
        TNLean.PEPS.EncodedFrame.SmallPatchRewrite.oldFrame,
        TNLean.PEPS.EncodedFrame.SmallPatchRewrite.newFrame,
        TNLean.PEPS.EncodedFrame.SmallPatchRewrite.Conditions,
        TNLean.PEPS.EncodedFrame.SmallPatchRewrite.AvoidsUntouched,
        TNLean.PEPS.EncodedFrame.Frame.innerHoles,
        TNLean.PEPS.EncodedFrame.Frame.outerHoles}
    \leanok
    \uses{def:peps_frame}
    Old and new one-sheet frames share their untouched holes, with the same
    encodings and tag owners; the remaining holes are the affected old holes and the affected new holes,
    $r_{\rm old}$ of the former. There are $m$ additional square patches
    $(c_i,u_i)$ with cylinder data, and a specified finite list of parties.
    Write $H^\mp(\mathcal F)$ for the unions of the physical samples of the
    inner and outer hole squares. The four conditions of
    \cite[Lemma~6.3]{OpenAI2026PolynomialPEPS}
    (\texttt{05-frames.tex}, lines 193--207) are:
    \begin{enumerate}
        \item every patch's outer square, and every affected old or new hole,
              avoids the physical outer footprints of the untouched holes;
        \item the union of the patch inner squares contains every site whose
              raw owner changes and the entire outer footprint of every
              affected old or new hole;
        \item in each patch's outer square, the sites outside
              $H^-(\mathcal F_{\rm old})$ have at most two old owners, and the
              sites outside $H^-(\mathcal F_{\rm new})$ have at most two new
              owners;
        \item all old and new raw owners on the patch footprints and the
              affected-hole footprints, and all affected tag owners, belong to
              the specified list.
    \end{enumerate}
    Patches may overlap, and affected old holes may overlap affected new
    holes. In the canonical coordinates used below, the tags of the untouched
    holes come first in both frames; Theorem~\ref{thm:peps_frame_small_rewrite}
    also covers arbitrary tag orderings.
\end{definition}

\begin{definition}[Canonical rewrite]
    \label{def:peps_frame_rewrite}
    \lean{TNLean.PEPS.EncodedFrame.SmallPatchRewrite.patchProj,
        TNLean.PEPS.EncodedFrame.SmallPatchRewrite.affectedRewrite,
        TNLean.PEPS.EncodedFrame.SmallPatchRewrite.layoutEquiv,
        TNLean.PEPS.EncodedFrame.SmallPatchRewrite.rewrite}
    \leanok
    \uses{def:peps_frame_rewrite_data}
    With $K_{\rm old}$ and $K_{\rm new}$ the encodings of the affected old
    and new holes alone and $P_1,\dots,P_m$ the patch projectors in a fixed
    order,
    \[
        M=K_{\rm new}P_m\cdots P_1K_{\rm old}^\dagger ,
    \]
    acting as the identity on the tags of the untouched holes. Its input
    layout is that of $\mathcal F_{\rm old}$ and its output layout that of
    $\mathcal F_{\rm new}$ (\texttt{05-frames.tex}, lines 222--232).
\end{definition}

\begin{lemma}[The rewrite is a contraction]
    \label{lem:peps_frame_rewrite_contraction}
    \lean{TNLean.PEPS.EncodedFrame.SmallPatchRewrite.norm_rewrite_le_one,
        Matrix.l2_opNorm_reindex_le,
        Matrix.l2_opNorm_kronecker_one_le}
    \leanok
    \uses{def:peps_frame_rewrite,lem:peps_frame_gram}
    $\lVert M\rVert\le1$.
\end{lemma}

\begin{proof}\leanok
    Every factor is a contraction. Tensoring with the identity on the
    untouched tags and relabelling coordinates do not increase the operator
    norm; the first estimate sums squared norms over the untouched tag
    slices and introduces no dimension factor.
\end{proof}

\begin{lemma}[Exact intertwining]
    \label{lem:peps_frame_rewrite_intertwining}
    \lean{TNLean.PEPS.EncodedFrame.SmallPatchRewrite.rewrite_mul_encoder}
    \leanok
    \uses{def:peps_frame_rewrite,lem:peps_frame_commute,lem:peps_frame_gram}
    Under condition (1),
    $MK_{\mathcal F_{\rm old}}=K_{\mathcal F_{\rm new}}P_m\cdots P_1
        Q_1\cdots Q_{r_{\rm old}}$, where $Q_1,\dots,Q_{r_{\rm old}}$ are the
    projectors of the affected old holes.
\end{lemma}

\begin{proof}\leanok
    On untouched tags $u$, affected old tags $b$ and affected new tags $a$,
    write $U_u$, $A_b$ and $B_a$ for the raw parts of the encodings. By
    condition (1) and disjointness of the old frame, $U_u$ commutes with
    $A_b^\dagger$, with every patch projector, with $B_a$ and with the
    affected old projectors. Summing over $b$,
    $\sum_bA_b^\dagger U_uA_b=U_u\sum_bA_b^\dagger A_b
        =U_uQ_1\cdots Q_{r_{\rm old}}$ by the frame Gram identity; then $U_u$
    commutes to the left of $B_aP_m\cdots P_1$.
\end{proof}

\begin{theorem}[Small-patch rewrite: canonical contraction]
    \label{thm:peps_frame_small_rewrite}
    \lean{TNLean.PEPS.EncodedFrame.SmallPatchRewrite.norm_rewrite_le_one_and_refVec_sub_le,
        TNLean.PEPS.EncodedFrame.SmallPatchRewrite.norm_rewrite_refVec_sub_le,
        TNLean.PEPS.EncodedFrame.exists_tagEquiv_of_perm,
        TNLean.PEPS.EncodedFrame.frameEncoder_submatrix_of_rawProd_eq,
        TNLean.PEPS.EncodedFrame.SmallPatchRewrite.exists_norm_le_one_and_refVec_sub_le_of_perm}
    \leanok
    \uses{lem:peps_frame_rewrite_contraction,lem:peps_frame_rewrite_intertwining,
        lem:peps_frame_telescoping}
    Under the four conditions, if every patch projector and every affected
    old-hole projector has error at most $\varepsilon$ on $\Omega$, the
    canonical map $M$ is a contraction and
    \[
        \lVert M\Omega_{\mathcal F_{\rm old}}-\Omega_{\mathcal F_{\rm new}}\rVert
        \le(m+r_{\rm old})\varepsilon .
    \]
    The same holds for frames whose tags are listed in any order, with $M$
    relabelled along the identification of the two tag orderings. This is the
    first assertion of \cite[Lemma~6.3]{OpenAI2026PolynomialPEPS}, with its
    hypotheses; the error estimate uses condition (1) only. No
    smallness of the errors of the untouched holes is assumed, and the
    projectors of overlapping patches need not commute.
\end{theorem}

\begin{proof}\leanok
    By the intertwining lemma,
    $M\Omega_{\mathcal F_{\rm old}}-\Omega_{\mathcal F_{\rm new}}
        =K_{\mathcal F_{\rm new}}(P_m\cdots P_1Q_1\cdots Q_{r_{\rm old}}-\id)\Omega$.
    The encoding is a contraction, and the telescoping lemma bounds the
    remaining vector by $(m+r_{\rm old})\varepsilon$. For another ordering of
    the tags, the raw parts of the encodings of holes with disjoint outer
    footprints commute, so moving every tag to the new position of its hole
    is a relabelling of the tag basis that carries one encoding to the other.
    Relabelling the coordinates of $M$ along it keeps $M$ a contraction and
    changes neither norm in the estimate.
\end{proof}

\begin{lemma}[Branch expansion]
    \label{lem:peps_frame_branches}
    \lean{TNLean.PEPS.EncodedFrame.PatchTags,
        TNLean.PEPS.EncodedFrame.cylProd,
        TNLean.PEPS.EncodedFrame.projProd_eq_sum_cylProd,
        TNLean.PEPS.EncodedFrame.SmallPatchRewrite.Branch,
        TNLean.PEPS.EncodedFrame.SmallPatchRewrite.branchRaw,
        TNLean.PEPS.EncodedFrame.SmallPatchRewrite.branchOp,
        TNLean.PEPS.EncodedFrame.SmallPatchRewrite.liftBranch,
        TNLean.PEPS.EncodedFrame.SmallPatchRewrite.affectedRewrite_eq_sum_branchOp,
        TNLean.PEPS.EncodedFrame.SmallPatchRewrite.rewrite_eq_sum_branch,
        TNLean.PEPS.EncodedFrame.SmallPatchRewrite.card_branch_le}
    \leanok
    \uses{def:peps_frame_rewrite}
    A complete branch $\beta=(a,s,b)$ chooses a tag $a$ of the affected new
    holes, one cylinder term of every patch, and a tag $b$ of the affected
    old holes. With $C_s$ the product of the chosen cylinder projectors,
    $M=\sum_\beta M_\beta$, where $M_\beta$ is
    $\ket a\bra b\otimes B_aC_sA_b^\dagger$ extended by the identity on the
    untouched tags. If every affected hole and every patch has at most $D$
    cylinder terms, there are at most $D^{\,r_{\rm new}+m+r_{\rm old}}$
    branches (\texttt{05-frames.tex}, lines 254--257).
\end{lemma}

\begin{proof}\leanok
    Expand every patch projector as the sum of its cylinder projectors and
    read off the entries of $K_{\rm new}P_m\cdots P_1K_{\rm old}^\dagger$ on
    the tag basis. The number of tag configurations of a list is the product
    of the numbers of cylinder terms of its members.
\end{proof}

\begin{lemma}[Wire classification of a branch]
    \label{lem:peps_frame_wires}
    \lean{TNLean.PEPS.EncodedFrame.SquarePatch.tagSample,
        TNLean.PEPS.EncodedFrame.patchSquares,
        TNLean.PEPS.EncodedFrame.holeSquares,
        TNLean.PEPS.EncodedFrame.SmallPatchRewrite.oldSample_subset_patchSquares,
        TNLean.PEPS.EncodedFrame.SmallPatchRewrite.newSample_subset_patchSquares,
        TNLean.PEPS.EncodedFrame.SmallPatchRewrite.exists_old_owners_of_patch_input,
        TNLean.PEPS.EncodedFrame.SmallPatchRewrite.exists_new_owners_of_patch_output,
        TNLean.PEPS.EncodedFrame.SmallPatchRewrite.ownerOld_eq_ownerNew_of_direct,
        TNLean.PEPS.EncodedFrame.SmallPatchRewrite.owners_mem_parties_of_mem_sample}
    \leanok
    \uses{def:peps_frame_rewrite_data}
    Assume the four conditions and fix a branch $(a,s,b)$. The selected
    square of a term $(j,\ell)$ is $D_j$, and that of the identity term of a
    square with empty outer sample is empty; the selected squares are read
    off from the terms of the branch position by position in the lists, so a
    patch listed twice may carry two different terms. Every site of a
    selected affected old-hole square, and of a
    selected affected new-hole square, lies in a selected patch square. The
    sites of a selected patch square not consumed by a selected old-hole
    square have at most two old owners, and those not received by a selected
    new-hole square have at most two new owners. A site in no selected patch
    square has the same old and new owner. All old and new owners of sites in
    selected squares belong to the specified list
    (\texttt{05-frames.tex}, lines 267--305).
\end{lemma}

\begin{proof}\leanok
    Every selected square contains its inner square and lies in its outer
    square (for an empty outer sample both are empty), so condition (2) gives the first statement and the statement on
    direct sites. A site of a patch outer square in
    $H^-(\mathcal F_{\rm old})$ is not in an untouched hole by condition (1),
    hence lies in the inner square of an affected old hole and so in its
    selected square; condition (3) bounds the owners of the remaining sites.
    The new case is symmetric, and condition (4) gives the last statement.
\end{proof}

\begin{lemma}[Rescaled branchwise approximation]
    \label{lem:peps_frame_rescale}
    \lean{TNLean.PEPS.EncodedFrame.norm_rescale_le,
        TNLean.PEPS.EncodedFrame.norm_rescale_sum_le_of_branches,
        TNLean.PEPS.EncodedFrame.truncationRank,
        TNLean.PEPS.EncodedFrame.truncationRank_spec}
    \leanok
    Let $M=\sum_\beta M_\beta$ be a contraction with $N$ branches and let
    $\lVert M_\beta-\widetilde M_\beta\rVert\le\eta$ for every branch. If
    $N\eta\le\tfrac12\delta$, then
    $M_\delta=(1+\tfrac12\delta)^{-1}\sum_\beta\widetilde M_\beta$ satisfies
    $\lVert M_\delta\rVert\le1$ and $\lVert M_\delta-M\rVert\le\delta$. For
    $N,g\ge0$ and $\delta>0$ there is an integer $k\ge1$ with
    $Ng/\sqrt k\le\tfrac12\delta$ (\texttt{05-frames.tex}, lines 331--341).
\end{lemma}

\begin{proof}\leanok
    The sum $\widetilde M$ is within $\tfrac12\delta$ of $M$, so
    $\lVert\widetilde M\rVert\le1+\tfrac12\delta$. Splitting
    $M_\delta-M=(1+\tfrac12\delta)^{-1}(\widetilde M-M)
        -\tfrac{\delta/2}{1+\delta/2}M$ bounds the error by
    $\delta/(1+\tfrac12\delta)\le\delta$. Take
    $k=\lceil(2Ng/\delta)^2\rceil+1$.
\end{proof}

\begin{lemma}[Reference error of an approximate rewrite]
    \label{lem:peps_frame_approx_reference}
    \lean{TNLean.PEPS.EncodedFrame.SmallPatchRewrite.norm_act_refVec_sub_le_of_approx,
        TNLean.PEPS.EncodedFrame.rewrite_error_le_inv_pow_twenty}
    \leanok
    \uses{thm:peps_frame_small_rewrite}
    Under condition (1), if $\lVert\Omega\rVert\le1$ and
    $\lVert M_a-M\rVert\le\delta$, then
    $\lVert M_a\Omega_{\mathcal F_{\rm old}}-\Omega_{\mathcal F_{\rm new}}\rVert
        \le\delta+(m+r_{\rm old})\varepsilon$. With $\varepsilon=L^{-60}$,
    $\delta=L^{-30}$ and $m+r_{\rm old}+1\le L^{10}$, this is at most
    $L^{-20}$ (\texttt{05-frames.tex}, lines 217--218 and 343--344).
\end{lemma}

\begin{proof}\leanok
    Split $M_a\Omega_{\mathcal F_{\rm old}}-\Omega_{\mathcal F_{\rm new}}$ as
    $(M_a-M)\Omega_{\mathcal F_{\rm old}}$ plus the error of the theorem, and
    use $\lVert\Omega_{\mathcal F_{\rm old}}\rVert\le1$. For the
    specialization, $L^{-60}\le L^{-30}$.
\end{proof}

\begin{lemma}[Chains of rank-one site operators as networks]
    \label{lem:peps_frame_chain_network}
    \lean{TNLean.PEPS.SiteChain.Step,
        TNLean.PEPS.SiteChain.chainOp,
        TNLean.PEPS.SiteChain.tensor,
        Matrix.ofFn_prod_apply,
        TNLean.PEPS.SiteChain.chainOp_apply}
    \leanok
    Let $L_0,\dots,L_{n-1}$ be operators
    $L_j=\ket{k_j}\bra{b_j}_{S_j}\otimes\id$ on the raw registers, where each
    $k_j$ and $b_j$ is a vector on the configurations of the finite set $S_j$
    or the product zero vector. Form the network whose vertices are the bras
    $\bra{b_j}$ and the kets $\ket{k_j}$, kept distinct, and whose internal
    edges join the output of a ket at a site to the input of the next bra
    containing that site, in the order $L_{n-1},\dots,L_0$ of application.
    The input of a bra at a site contained in no earlier factor is an open
    input leg, and the output of a ket at a site contained in no later factor
    is an open output leg, except that the legs of a product zero vector are
    fixed to $0$. Then the entry of $L_0\cdots L_{n-1}$ at output
    configuration $\sigma$ and input configuration $\tau$ vanishes unless
    $\sigma$ and $\tau$ agree at the sites contained in no factor and carry
    $0$ at the fixed legs, and otherwise equals the contraction of the network
    evaluated at the open legs read off from $\sigma$ and $\tau$. No edge
    joins a vertex to itself (\texttt{05-frames.tex}, lines 274--281).
\end{lemma}

\begin{proof}\leanok
    Expand the product as a sum over paths of intermediate configurations.
    A path of nonzero weight keeps the value of a site across factors that do
    not contain it, so it is determined by the values on the internal edges,
    and conversely every labelling of the internal edges determines such a
    path; along it the weight is the product of the vertex tensors.
\end{proof}

\begin{lemma}[Whole-group truncation of a chain]
    \label{lem:peps_frame_chain_truncation}
    \lean{Matrix.l2_opNorm_le_of_blocks,
        TNLean.PEPS.SiteChain.vertexNormSq_tensor_le_one,
        TNLean.PEPS.SiteChain.approxOp,
        TNLean.PEPS.SiteChain.norm_approxOp_le,
        TNLean.PEPS.SiteChain.termOp,
        TNLean.PEPS.SiteChain.exists_chainOp_truncation}
    \leanok
    \uses{lem:peps_frame_chain_network,thm:peps_group_tensor}
    In the setting of the previous lemma, let every $k_j$ and $b_j$ have norm
    at most one, and let $G$ be a set of vertices containing every vertex
    with an open leg. For every $k\ge1$ there are orthonormal families of at
    most $k$ vectors on the open-leg groups and coefficients of modulus at
    most one such that $L_0\cdots L_{n-1}$ differs in operator norm by at most
    $|G|/\sqrt k$ from the sum of the corresponding product terms, of which
    there are at most $k^{|G|}$. Each product term is the tensor product of
    one normalized vector on the open output legs of each ket, one
    normalized covector on the open input legs of each bra, the product zero
    vectors at the fixed legs, and the identity on the sites contained in no
    factor. No constant depends on the dimensions of the sites
    (\texttt{05-frames.tex}, lines 306--329).
\end{lemma}

\begin{proof}\leanok
    The vertex tensors are normalized, so the whole-group truncation lemma
    applies to the contraction $Z$. An operator read off from a tensor $W$ on
    the open legs is $W$, reshaped as a matrix from inputs to outputs,
    tensored with the identity on the remaining sites, so its operator norm is
    at most the largest Hilbert--Schmidt norm of its diagonal blocks, hence
    at most $\lVert W\rVert$. Apply this to $W=Z-Z_k$.
\end{proof}

\begin{lemma}[Network of a branch]
    \label{lem:peps_frame_branch_network}
    \lean{TNLean.PEPS.EncodedFrame.SquarePatch.tagVec,
        TNLean.PEPS.EncodedFrame.SmallPatchRewrite.branchSteps,
        TNLean.PEPS.EncodedFrame.SmallPatchRewrite.branchRaw_eq_chainOp,
        TNLean.PEPS.EncodedFrame.SmallPatchRewrite.patchVertices,
        TNLean.PEPS.EncodedFrame.SmallPatchRewrite.card_patchVertices_le,
        TNLean.PEPS.EncodedFrame.SmallPatchRewrite.mem_patchVertices_of_isOpen,
        TNLean.PEPS.EncodedFrame.SmallPatchRewrite.exists_owners_of_isOpen_inl,
        TNLean.PEPS.EncodedFrame.SmallPatchRewrite.exists_owners_of_isOpen_inr}
    \leanok
    \uses{lem:peps_frame_branches,lem:peps_frame_wires,
        lem:peps_frame_chain_network}
    Assume the four conditions and fix a branch $(a,s,b)$. Its raw operator
    $B_aC_sA_b^\dagger$ is the ordered product of the factors
    $\ket0\bra{u'_r}$ of the affected new holes, $\ket{w_i}\bra{w_i}$ of the
    patch terms, and $\ket{u_t}\bra0$ of the affected old holes. In its
    network only the bras and kets of the patch terms have open legs, so at
    most $2m$ vertices do; the open input legs of a patch bra have at most
    two old owners, the open output legs of a patch ket have at most two new
    owners, and all these owners belong to the specified list
    (\texttt{05-frames.tex}, lines 257--305).
\end{lemma}

\begin{proof}\leanok
    The factorization follows from the definitions of the hole encodings and
    the cylinder projectors. A decoded old-hole ket meets a later patch bra
    at every site of its square, and a new-hole bra meets an earlier patch
    ket at every site of its square, by the wire classification; the zero
    vectors carry no open legs. An open input of a patch bra is consumed by
    no selected old-hole square, so the wire classification bounds its old
    owners; the output case is symmetric.
\end{proof}

\begin{theorem}[Small-patch rewrite: approximation by product terms]
    \label{thm:peps_frame_small_rewrite_product_terms}
    \lean{TNLean.PEPS.EncodedFrame.SmallPatchRewrite.exists_branch_truncation,
        TNLean.PEPS.EncodedFrame.SmallPatchRewrite.liftBranchLinearMap,
        TNLean.PEPS.EncodedFrame.SmallPatchRewrite.liftBranchLinearMap_apply,
        TNLean.PEPS.EncodedFrame.SmallPatchRewrite.liftBranch_sub,
        TNLean.PEPS.EncodedFrame.SmallPatchRewrite.norm_liftBranch_le,
        TNLean.PEPS.EncodedFrame.SmallPatchRewrite.exists_contractive_approx}
    \leanok
    \uses{lem:peps_frame_branch_network,lem:peps_frame_chain_truncation,
        lem:peps_frame_rescale}
    Assume the four conditions, and let $N$ be the number of branches and
    $k=\lceil(4Nm/\delta)^2\rceil+1$ for $\delta>0$, for instance
    $\delta=L^{-a}$. Every branch is approximated in operator norm within
    $2m/\sqrt k$ by at most $k^{2m}$ product terms with coefficients of
    modulus at most one, and the rescaled sum
    $M_\delta=(1+\tfrac12\delta)^{-1}\sum_\beta\ket a\bra b\otimes
        \bigl(\sum_ic_{\beta,i}T_{\beta,i}\bigr)$, extended by the identity on
    the untouched tags, is a contraction with $\lVert M_\delta-M\rVert\le
        \delta$. Each product term $T_{\beta,i}$ is the tensor product of
    normalized vectors on the open output groups of the patch kets, normalized
    covectors on the open input groups of the patch bras, product zero
    vectors on the selected hole squares, and identities; by
    Lemma~\ref{lem:peps_frame_branch_network} each open group lies on at most
    two parties of the specified list. This formalizes the branch-network,
    truncation and rescaling steps of the proof of the second assertion of
    \cite[Lemma~6.3]{OpenAI2026PolynomialPEPS}. The reading of the product
    terms as allowed monomials is
    Theorem~\ref{thm:peps_frame_small_rewrite_monomials}, and the polynomial
    bound on the number of terms is
    Theorem~\ref{thm:peps_frame_small_rewrite_polynomial}.
\end{theorem}

\begin{proof}\leanok
    Apply the truncation of a chain to the network of each branch with the
    patch vertices as groups, and assemble with the rescaling lemma, using
    $N\cdot 2m/\sqrt k\le\tfrac12\delta$.
\end{proof}

\begin{theorem}[Small-patch rewrite: polynomial approximation]
    \label{thm:peps_frame_small_rewrite_polynomial}
    \lean{TNLean.PEPS.EncodedFrame.SmallPatchRewrite.exists_allowed_monomial_approx_polynomial}
    \leanok
    \uses{thm:peps_frame_small_rewrite,lem:peps_frame_branches,
        lem:peps_frame_wires,lem:peps_frame_rescale,
        lem:peps_frame_approx_reference,thm:peps_group_tensor,
        thm:peps_frame_small_rewrite_product_terms,
        thm:peps_frame_small_rewrite_monomials}
    In the setting of the frame construction, where the projector of every
    hole and every additional patch is a sum of $\sum_jd_j\le CL^c$
    cylinder projectors for fixed $C$ and $c$ (\texttt{05-frames.tex}, lines
    34--46), assume the four conditions with boundedly many affected holes and
    additional patches. Then for every fixed $a>0$ the canonical map $M$
    admits a contractive approximation $M_a$ with $\lVert M_a-M\rVert\le
        L^{-a}$ and a polynomial sum expansion into monomials allowed by
    \cite[Theorem~5.2]{OpenAI2026PolynomialPEPS}, using only the specified
    parties. In particular the rewrite is a bounded change with
    reference-vector error $O(L^{-20})$. This is the second assertion of
    \cite[Lemma~6.3]{OpenAI2026PolynomialPEPS}.
\end{theorem}

\begin{proof}\leanok
    Expand $M$ into its polynomially many branches. In a branch
    $\beta=(a,s,b)$, let the old affected holes carry the terms with selected
    squares $D^{\rm old}_t$ and unit vectors $u_t$, the new affected holes
    the terms with selected squares $D^{\rm new}_r$ and unit vectors $u'_r$,
    and the patches the terms with selected squares $E_i$ and unit vectors
    $w_i$. Then $M_\beta=\ket a\bra b\otimes B_aC_sA_b^\dagger$ with
    \[
        A_b^\dagger=\prod_t\bigl(\ket{u_t}\bra0^{\otimes D^{\rm old}_t}
        \bigr)_{D^{\rm old}_t}\otimes\id,\quad
        C_s=\prod_i\bigl(\ket{w_i}\bra{w_i}\bigr)_{E_i}\otimes\id,\quad
        B_a=\prod_r\bigl(\ket0^{\otimes D^{\rm new}_r}\bra{u'_r}
        \bigr)_{D^{\rm new}_r}\otimes\id,
    \]
    the factors of $C_s$ taken in the order of the patches. Factor out the
    tag ket $\ket a$ and bra $\bra b$, the zero bras on the squares
    $D^{\rm old}_t$, the zero kets on the squares $D^{\rm new}_r$, and the
    identity on the sites outside $\bigcup_iE_i$, whose old and new owners
    agree. The residual network $Z$ has the vertices $\ket{u_t}$,
    $\bra{w_i}$, $\ket{w_i}$ and $\bra{u'_r}$, the bra and the ket of each
    patch term being distinct vertices; each site wire joins a ket to the
    next bra containing that site, so there are no self-contractions. By the
    wire classification, $\ket{u_t}$ has no open output legs and
    $\bra{u'_r}$ has no open input legs. The open legs of $Z$ are the input
    legs of $\bra{w_i}$ at sites of $E_i$ outside $\bigcup_tD^{\rm old}_t$
    not reached by an earlier ket, which have at most two old owners, and
    the output legs of $\ket{w_i}$ at sites of $E_i$ outside
    $\bigcup_rD^{\rm new}_r$ not met by a later bra, which have at most two
    new owners. The whole-group truncation
    lemma (Theorem~\ref{thm:peps_group_tensor}) approximates its tensor
    within $g/\sqrt k$ by at most $k^g$ product
    terms with coefficients of modulus at most one, each an allowed
    monomial, and the Hilbert-space error bounds the operator error of the
    branch. The rescaling lemma assembles $M_a$, and the reference estimate
    with $a=30$ gives the final assertion. The branch network, its truncation
    and the assembly are formalized in the approximation by product terms
    (Theorem~\ref{thm:peps_frame_small_rewrite_product_terms}), and the
    reading of each product term as an allowed monomial on the registers of
    the frames, with at most $Nk^{2m}$ monomials, in
    Theorem~\ref{thm:peps_frame_small_rewrite_monomials}. With
    $\sum_jd_j\le CL^c$ and $m,r_{\rm old},r_{\rm new}\le b$, the number of
    branches is at most $\bigl((C+1)L^c\bigr)^{3b}$, and for $\delta=L^{-a}$,
    $L\ge1$ the count is at most
    $(C+1)^{3b}\bigl(16b^2(C+1)^{6b}+1\bigr)^{2b}L^{3bc+4b(3bc+\lceil
                a\rceil)}$, with at most $2b$ normalized pair sources and $2b$ normalized
    pair effects in each monomial \cite{gap:polypeps_small_rewrite_monomials}.
\end{proof}
